% !PRINTED 498
\weberpage{498}

Setzen wir
\weq{3}{a_0 = m\,e_0,\quad a_1 = m\,e_1,\quad a_2 = m\,e_2 \ldots,}
so sind die $e_0, e_1, e_2, \ldots$ ganze Zahlen ohne gemeinsamen Theiler;
\weq{4}{E(x,y,z,\ldots) = \sum e\,x^r y^s z^t \ldots}
ist eine primitive Function und es wird
\[ \Phi = m\,E. \]

Diese Functionen $\Phi$ haben wir im \S.\,2 des ersten Bandes betrachtet und haben dort von ihnen den Satz bewiesen:

3. \gesperrt{Der Theiler eines Productes von zwei oder mehr ganzen Functionen $\Phi$ ist gleich dem Product der Theiler der einzelnen Factoren.}

Die ganzen Functionen in $\Omega$ h\"angen ausser von den Variablen von einer algebraischen Zahl $\Theta$ ab, enthalten aber sonst nur rationale Zahlenco\"efficienten. Bezeichnen wir eine solche Function mit $\varphi(\Theta, x, y, z, \ldots)$, so erhalten wir die \gesperrt{conjugirten Functionen} $\varphi, \varphi_1, \varphi_2, \ldots$ oder
\weq{5}{\varphi(\Theta,x,y,z,\ldots),\quad \varphi(\Theta_1,x,y,z,\ldots),\quad \varphi(\Theta_2,x,y,z,\ldots),\ \ldots}
wenn wir f\"ur $\Theta$ die s\"ammtlichen Wurzeln der irreduciblen Gleichung $f(x)=0$ [\S.\,132, (2)] einsetzen. Diese conjugirten Functionen k\"onnen auch zum Theil einander gleich sein.

Sie sind alle einander gleich, wenn $\varphi$ eine Function in $R$ ist, und es ist umgekehrt $\varphi$ eine Function in $R$, wenn die conjugirten Functionen alle einander gleich sind; denn es ist dann, wenn wir unter $S(\varphi)$ die Summe der conjugirten Functionen (die Spur) verstehen,
\[ n\varphi = S(\varphi), \]
und $S(\varphi)$ ist eine ganze Function, deren Co\"efficienten symmetrische Functionen der $n$ Wurzeln $\Theta$, d. h. rationale Zahlen, sind.

Zu den ganzen Functionen in $R$ geh\"ort auch die \gesperrt{Norm von} $\varphi$, d. h. das Product
\weq{6}{N(\varphi) = \varphi\,\varphi_1\,\varphi_2 \ldots,}
denn alle Co\"efficienten dieser Function sind symmetrische Functionen der $\Theta$. Diese Function ist theilbar durch $\varphi$, und wenn wir
\[ N(\varphi) = \varphi\,\varphi' \]
setzen, so sind sowohl $\varphi$ als $\varphi'$ ganze Functionen in $\Omega$. Denn $\varphi'$ ist als Product $\varphi_1\varphi_2\ldots$ eine ganze Function, und die Co\"effi-
%CONT>>
